\chapter{Mononucleosis (Mono) Testing}

\section*{What Is This Test?}
Mononucleosis testing detects acute Epstein-Barr virus (EBV) infection, the cause of infectious mononucleosis ("mono"). EBV is a ubiquitous gamma-herpesvirus (HHV-4) transmitted via saliva ("kissing disease"), with >90\% of adults worldwide seropositive \cite{cohen2000epstein}. Primary EBV infection in early childhood is usually asymptomatic or causes mild, nonspecific illness \cite{dunmire2018primary}. In adolescents and young adults (ages 15--24), primary EBV infection causes the classic mononucleosis syndrome: fever (38--39\textdegree{}C), severe exudative pharyngitis (often with thick, white tonsillar exudates resembling streptococcal pharyngitis or diphtheria), bilateral posterior cervical lymphadenopathy, fatigue (profound, lasting weeks to months), and splenomegaly (50--60\% on physical exam, nearly universal on ultrasound) \cite{ebell2004epstein}. Laboratory hallmarks include atypical lymphocytosis (>10\% of total lymphocytes on peripheral smear; >50\% in some cases) and mildly elevated transaminases (ALT, AST, 2--3x upper limit). Splenic rupture is a rare (0.1--0.5\%) but potentially fatal complication occurring most commonly in the second to third week of illness, often associated with trauma \cite{luzuriaga2010infectious}. Testing includes: heterophile antibody test (Monospot), EBV-specific serology (viral capsid antigen/VCA IgG and IgM, early antigen/EA IgG, Epstein-Barr nuclear antigen/EBNA IgG), and EBV PCR.

\section*{Alternative Names}
Mono Test; Monospot; Heterophile Antibody Test; EBV Serology; EBV Antibody Panel; Infectious Mononucleosis Testing; VCA IgG/IgM; EBNA IgG; Paul-Bunnell Test (historical)

\section*{Principle}
Heterophile antibody test (Monospot): A latex agglutination or immunochromatographic assay detecting heterophile IgM antibodies that agglutinate horse, sheep, or bovine erythrocytes. These antibodies are not specific to EBV; they are produced from polyclonal B-cell activation during acute EBV infection. Monospot sensitivity is 70--90\% in immunocompetent adolescents/adults, but only 25--50\% in children under 4 years, in whom heterophile antibodies are often absent. Specificity is 95--98\%. False positives occur in: autoimmune disease, lymphoma, leukemia, hepatitis, malaria, rubella, and HIV seroconversion. Results available in 5--10 minutes. EBV-specific serology: CLIA/EIA detecting IgG and IgM antibodies specific to EBV antigens: VCA IgM (appears at symptom onset; marker of acute primary infection; declines over 4--8 weeks, though may persist for months), VCA IgG (appears within 2--4 weeks of symptoms; peaks at 2--3 months; persists for life), EA IgG (appears during acute infection; marker of active viral replication; declines to undetectable in most within 3--6 months; may persist in some healthy individuals and is not a standalone marker of chronic EBV), EBNA IgG (appears 6--12 weeks after primary infection---late; persists for life; the presence of EBNA IgG in the setting of acute symptoms effectively excludes primary EBV infection, as EBNA develops late). The pattern of reactivity differentiates acute primary infection from past infection and from reactivation: Acute primary: VCA IgM +, VCA IgG +/-, EA IgG +/-, EBNA IgG -. Past infection: VCA IgM -, VCA IgG +, EA IgG - (or low +), EBNA IgG +. Reactivation (immunocompromised): VCA IgM -/+, VCA IgG +++ (very high), EA IgG ++, EBNA IgG + \cite{klutts2009evidence}. EBV PCR: Quantitative real-time PCR for EBV DNA in whole blood, plasma, or CSF. Used primarily for monitoring EBV-associated lymphoproliferative disease in transplant recipients (post-transplant lymphoproliferative disorder/PTLD) and for diagnosis of EBV-associated malignancies (nasopharyngeal carcinoma, lymphoma). Not used for routine diagnosis of infectious mononucleosis.

\section*{History}
Infectious mononucleosis was described as a clinical entity ("glandular fever" or "Drüsenfieber") by Emil Pfeiffer in 1889. The Paul-Bunnell test (heterophile antibody detection by sheep erythrocyte agglutination) was described in 1932. EBV was discovered by Michael Epstein and Yvonne Barr in 1964 from cultured Burkitt lymphoma cells. Werner and Gertrude Henle demonstrated that EBV is the cause of infectious mononucleosis in 1968. EBV was the first virus shown to cause cancer in humans (Burkitt lymphoma, nasopharyngeal carcinoma, Hodgkin lymphoma, post-transplant lymphoproliferative disorder). The Monospot rapid latex agglutination test was developed in the 1970s and remains the most widely used point-of-care test. EBV-specific serology (VCA, EA, EBNA) was developed in the 1980s and refined with ELISA technology.

\section*{Why This Test Is Ordered}
\begin{itemize}
  \item Diagnosis of infectious mononucleosis in adolescents and young adults presenting with the classic triad: fever, severe exudative pharyngitis (often with thick tonsillar exudates refractory to antibiotics), and bilateral posterior cervical lymphadenopathy, plus fatigue, splenomegaly, and atypical lymphocytosis
  \item Differentiation of acute EBV infection from group A streptococcal pharyngitis, which has overlapping clinical features but requires antibiotic therapy. Clinicians should note that ampicillin or amoxicillin administered during acute EBV infection causes a characteristic diffuse maculopapular rash in >90\% of patients (this is not a true penicillin allergy)
  \item Evaluation of prolonged fatigue (\textgreater 6 months) after an acute mono-like illness to confirm antecedent EBV infection (EBV serology documenting past infection)
  \item Evaluation of lymphadenopathy with atypical features, splenomegaly, or unexplained atypical lymphocytosis on CBC with differential
  \item When the Monospot is negative but clinical suspicion remains high (particularly in children under 4, early in the illness, or late presentation), EBV-specific serology should be performed
  \item Evaluation of suspected chronic active EBV (CAEBV), a rare, severe progressive disease seen primarily in East Asian and Latin American populations, characterized by persistent or recurrent IM-like symptoms, abnormal EBV serology (markedly elevated VCA IgG and EA IgG), and high EBV DNA in peripheral blood and/or tissue
  \item Monitoring of EBV viral load in transplant recipients at risk for PTLD
\end{itemize}

\section*{Primary vs Secondary Test}
The Monospot (heterophile antibody test) is the primary screening test for suspected infectious mononucleosis in immunocompetent adolescents and adults, due to its speed, low cost, and acceptable sensitivity. EBV-specific serology is the primary test when the Monospot is negative with high clinical suspicion, in children under 4, and when distinguishing acute from past infection is necessary. EBV PCR is a secondary test for routine mononucleosis diagnosis but primary for PTLD surveillance.

\section*{How This Test Is Performed}
Monospot: Venous blood (SST tube) or fingerstick blood is sufficient. Test kits are lateral flow or latex agglutination-based. Results read in 5--10 minutes. EBV serology: Venous blood in SST tube; serum separated; 5--7 mL yields sufficient volume for the full panel (VCA IgM, VCA IgG, EA IgG, EBNA IgG). EBV PCR: EDTA whole blood or plasma; plasma is preferred for transplant recipient monitoring.

\section*{What Preparation Is Needed}
No special preparation is required. Patients do not need to fast. If the patient has received heterophile antibodies passively (e.g., from IVIG, blood transfusion, or heterophile antibody preparations), the Monospot may be falsely positive.

\section*{Contraindications and Precautions}
The Monospot has significant limitations: (1) Reduced sensitivity in children under 4 years (25--50\% false-negative rate). (2) False-negative in the first week of illness (heterophile antibodies may not have risen to detectable levels); if clinical suspicion is high, repeat the Monospot in 5--7 days or order EBV-specific serology. (3) False-positive results in autoimmune disease (SLE, rheumatoid arthritis), lymphoma, leukemia, viral hepatitis, malaria, rubella, HIV seroconversion, and pancreatic cancer. A positive Monospot in an atypical clinical context should prompt consideration of these alternate diagnoses. (4) Heterophile antibodies may persist for up to 12 months after acute infection; a positive Monospot in a patient currently ill but with a history of mono 3--6 months ago does not necessarily represent a current EBV infection. EBV serology is specific and is the gold standard. Anti-VCA IgM can cross-react with CMV IgM in some assays. EBNA IgG may be absent in immunocompromised patients despite past EBV infection. EA IgG is elevated in 10--20\% of healthy EBV-seropositive individuals and is not a standalone marker of EBV reactivation or chronic EBV disease. Splenic palpation should be performed gently, if at all, given the risk of splenic rupture.

\section*{Risks}
Standard venipuncture risks. The primary risk is splenic rupture, which is not from testing but from the disease. Patients with suspected or confirmed mononucleosis should avoid contact sports, heavy lifting, and abdominal trauma for a minimum of 3--4 weeks after symptom onset (and until splenomegaly has resolved clinically or radiographically). Ultrasonography to document splenic size can guide return-to-play decisions for athletes.

\section*{What Happens After the Test}
Monospot: results in 5--10 minutes. EBV serology: results in 1--24 hours. EBV PCR: results in 24--72 hours. Management of infectious mononucleosis is supportive: rest, adequate hydration, acetaminophen or NSAIDs for fever and throat pain. Corticosteroids (prednisone 40--60 mg/day tapered over 5--7 days) are indicated for: impending airway obstruction (tonsillar hypertrophy), severe thrombocytopenia, hemolytic anemia, or severe constitutional symptoms preventing oral intake. Antibiotics (amoxicillin, ampicillin) should be avoided because of the near-universal maculopapular rash they induce during EBV infection. Acyclovir and other antivirals have no clinically significant benefit in acute mononucleosis and are not recommended. Splenomegaly should be monitored; full return to contact sports is generally permitted at 4--6 weeks if the spleen is nonpalpable and ultrasound confirms normal size.

\section*{Understanding Results}

\subsection*{Below Normal}
\begin{itemize}
  \item \textbf{Monospot negative:} Heterophile antibodies not detected. May indicate: (a) absence of EBV infection; (b) early acute infection (first 5--7 days) before heterophile antibodies have developed; (c) EBV infection in a child under 4 (25--50\% sensitivity, often heterophile-negative); (d) CMV mononucleosis (which is heterophile-negative); or (e) toxoplasmosis, acute HIV, or HHV-6, all of which can cause mononucleosis-like illness
  \item \textbf{EBV serology: VCA IgM negative, VCA IgG negative, EBNA IgG negative:} EBV seronegative (no prior EBV infection). The patient is susceptible to primary EBV infection. The current illness is NOT due to EBV. Differential includes CMV, HIV, toxoplasmosis, HHV-6, streptococcal pharyngitis, and other causes of pharyngitis
  \item \textbf{EBV serology: VCA IgG positive, EBNA IgG positive:} Past EBV infection (latent). The current illness is NOT acute EBV. These results are typical of >90\% of adults worldwide who have been infected with EBV in the past
\end{itemize}

\subsection*{Above Normal}
\begin{itemize}
  \item \textbf{Monospot positive:} Heterophile antibodies detected. In a patient with compatible clinical syndrome, diagnostic of infectious mononucleosis. Management is supportive. False-positive rate is 2--5\%; if the clinical picture is atypical, consider alternative causes of a positive Monospot (HIV, autoimmune disease, lymphoma). The Monospot may remain positive for up to 12 months after acute infection
  \item \textbf{EBV serology: VCA IgM positive, VCA IgG positive or negative, EBNA IgG negative:} Acute primary EBV infection (infectious mononucleosis). This pattern is diagnostic. The absence of EBNA IgG confirms the infection is recent (<6--12 weeks). The patient should avoid contact sports and abdominal trauma for 4--6 weeks
  \item \textbf{EBV serology: VCA IgG markedly elevated, EA IgG elevated, EBNA IgG positive:} EBV reactivation. This pattern may be seen in immunocompromised patients (transplant, HIV), during physiologic stress, or occasionally during other acute febrile illnesses. CAEBV should be considered if this pattern is accompanied by persistent severe symptoms, organomegaly, and high EBV DNA loads
  \item \textbf{Atypical lymphocytosis (>10\% atypical lymphocytes on peripheral smear):} Hallmark of infectious mononucleosis. The "atypical" lymphocytes are activated CD8+ cytotoxic T cells responding to EBV-infected B cells. CBC with differential is a useful adjunct: >50\% lymphocytes with >10\% atypical forms is highly suggestive of EBV
  \item \textbf{EBV PCR (viral load) elevated:} Used primarily in the setting of PTLD, CAEBV, or to monitor EBV-associated lymphomas. In acute infectious mononucleosis, EBV DNA in blood is detectable but generally moderate (10$^{2}$--10$^{4}$ copies/mL)
\end{itemize}

\section*{Normal Results}
Monospot: \textbf{Negative (no heterophile antibodies detected)}. EBV serology: \textbf{VCA IgG positive, VCA IgM negative, EBNA IgG positive} (evidence of past/latent EBV infection; most adults). EBV-naive: \textbf{VCA IgG negative, VCA IgM negative, EBNA IgG negative}. EBV PCR: \textbf{Not detected} or \textbf{below LLOQ}.

\section*{Follow-up Tests}
\begin{itemize}
  \item \textbf{EBV-specific serology (VCA IgM, VCA IgG, EA IgG, EBNA IgG):} When Monospot is negative with high clinical suspicion, in children under 4, or for definitive staging of EBV infection
  \item \textbf{CMV serology (CMV IgM/IgG) and CMV PCR:} When the Monospot is negative and the clinical picture resembles mononucleosis (heterophile-negative mononucleosis); CMV is the most common cause of heterophile-negative mononucleosis syndrome
  \item \textbf{HIV testing (4th-generation antigen/antibody + HIV-1 RNA if indicated):} Acute HIV seroconversion can mimic mononucleosis (fever, pharyngitis, lymphadenopathy, atypical lymphocytosis) and is a critical diagnosis to recognize
  \item \textbf{Toxoplasma serology:} Acute toxoplasmosis can cause a mononucleosis-like syndrome with posterior cervical lymphadenopathy (Piringer-Kuchinka syndrome), especially in individuals with exposure to cats or undercooked meat
  \item \textbf{CBC with differential and peripheral blood smear:} Evaluates atypical lymphocytosis, which supports EBV. Also evaluates for thrombocytopenia, hemolytic anemia, or lymphoma (if smear shows blasts or abnormal lymphocytes)
  \item \textbf{Abdominal ultrasound:} To assess splenomegaly and guide return-to-play decisions for athletes with confirmed mononucleosis
  \item \textbf{Throat culture or rapid strep test:} To exclude concurrent group A streptococcal pharyngitis, which can co-exist with EBV and requires antibiotic therapy
\end{itemize}

\section*{References}
\bibliographystyle{unsrtnat}
\bibliography{references}

\clearpage
